\begin{exo}{}
	On considère la figure ci-dessous composée d'un carré $ABCD$ de côté $1$ et de deux triangles équilatéraux $DIC$ et $BJC$.

	\begin{center}
		\begin{tikzpicture}[scale=1.6, >=Stealth]
			\coordinate (I) at (1,1.732);
			\coordinate (J) at (3.732,1);
			\draw[dashed, gray] (0,2) -- (J);
			\draw[thick] (0,2) node {$\bullet$} node[above left] {$A$} -- (2,2) node {$\bullet$} node[above right] {$B$} -- (2,0) node {$\bullet$} node[below right] {$C$} -- (0,0) node {$\bullet$} node[below left] {$D$} -- cycle;
			\draw[thick] (0,0) -- (I) node {$\bullet$} node[above] {$I$} -- (2,0);
			\draw[thick] (2,2) -- (J) node {$\bullet$} node[right] {$J$} -- (2,0);
			\draw[->, very thick] (0,0) -- (0,2);
			\draw[->, very thick] (0,0) -- (2,0);
		\end{tikzpicture}
	\end{center}

	On rappelle que dans un triangle équilatéral de côté $a$, toutes ses hauteurs ont pour longueur $\dfrac{a\sqrt{3}}{2}$.

	On munit le plan du repère orthonormé $(D\;; \vec{DC}\;, \vec{DA})$.

	\begin{enumerate}
		\item % Partie 1 : Coordonnées des sommets du carré
		      Donner, par lecture graphique, les coordonnées des points $D$, $C$, $A$ et $B$.

		\item % Partie 2 : Hauteur et coordonnées du point I
		      On s'intéresse au point $I$.
		      \begin{enumerate}
			      \item Justifier que la hauteur issue de $I$ dans le triangle équilatéral $DIC$ mesure $\dfrac{\sqrt{3}}{2}$.
			      \item En déduire les coordonnées exactes de $I$.
		      \end{enumerate}

		\item % Partie 3 : Hauteur et coordonnées du point J
		      On s'intéresse au point $J$.
		      \begin{enumerate}
			      \item Déterminer la distance de $J$ à la droite $(BC)$.
			      \item En déduire les coordonnées exactes de $J$.
		      \end{enumerate}

		\item % Partie 4 : Alignement
		      On souhaite démontrer que les points $A$, $I$ et $J$ sont alignés.
		      \begin{enumerate}
			      \item Déterminer les coordonnées de $\vec{AI}$ et $\vec{AJ}$.
			      \item Établir que les vecteurs $\vec{AI}$ et $\vec{AJ}$ sont colinéaires. \item Conclure quant à l'alignement de $A$, $I$ et $J$.
		      \end{enumerate}
	\end{enumerate}
\end{exo}
